% !TEX root = userguide.tex

\def\headdate{28th January 2014}

\section{Viscoelastic fluid flow using DG (explicit stress formulation)}

In this section some examples using the equations for viscoelastic fluid flow
will be given after the theory in the next section.

\subsection{Weak form of the viscoelastic fluid equations}
\label{sec:DGweak}

We assume the user is already familiar with Sec.~\ref{sec:weakviscoelastic}.
We will only explain the differences between SUPG and DG here. 

The equations for the flow of a single-mode viscoelastic fluid
in a region $\Omega$ are given by Eqs.~\eqref{eq:ve1}--\eqref{eq:ve3}.

For the discontinuous Galerkin \fem\ we first start with a finite element mesh
before defining the weak form:
\[
\Omega=\cup_e\Omega_e, \quad\Omega_{e}\cap\Omega_{f}=\emptyset\quad
   \text{for }e\neq f
\]
The weak form can be now be stated, element wise for the constitutive equation,
as:
Find $\vek u$, $p$ and $\ten c$ such that
\begin{alignat}{2}
 \big(\vek v,\rho(\pderiv{\vek u}t+\vek u\cdot\nabla\vek u)\big) +
  \big((\nabla\vek v)^T,2\eta_s\ten D\big)
    &-(\nabla\cdot\vek v,p)+\big((\nabla\vek v)^T,\ten\tau(\ten c)\big)\notag\\
    &=(\vek v,\vek t_{\text{N}})_{\Gamma_{\text{N}}}+(\vek v,\rho\vek b)
&\qquad&\text{for all }\vek v\\
  \,(q,\nabla\cdot\vek u)&=0,&\qquad&\text{for all }q\\
 \int_{\Omega_e}\ten d:\big(\pderiv{\ten c}{t}
+\vek u\cdot\nabla\ten c-\ten g(\ten L,\ten c)\big)\D x
    &+\int_{\Gamma^e_{\text{in}}}(\vek u\cdot\vek n)\ten d:[\ten c]\D s = 0,
  &\qquad&\text{for all }\ten d,\Omega_e
   \label{eq:weakcDG}
\end{alignat}
using appropriate spaces for $\vek u$, $p$, $\ten c$, $\vek v$, $q$ and $\ten
d$.
The function space for $\ten c$ and $\ten d$ is discontinuous on the
element boundaries.
In equation~\eqref{eq:weakcDG} $\Gamma^e_{\text{in}}$ is the
inflow boundary of element $e$ and
$[\ten c]=\ten c^+-\ten c$ is the ``jump'' across that element boundary.
Also, $\vek n$ is the outwardly directed unit normal vector on the boundary of
the element and $\ten c^+$ is
the `outside' value of $\ten c$.
 The function $\ten g(\ten L,\ten c)$ is given by
\begin{equation}
\ten g(\ten L,\ten c)=\ten L\cdot\ten c+\ten c\cdot\ten L^T
 -\ten f(\ten c)
\end{equation}

The weak form can be used to obtain an approximate solution
by interpolating $\vek u$, $p$ and $\ten c$ 
(and similarly $\vek v$, $q$ and $\ten d$) by
polynomials on each element:
\begin{align}
    \vek u_h&=\tsum_k \phi_k(\vek x)\vek u_k=\col\phi^T(\vek x)\col{\vek u}\\
    p_h&=\tsum_k \psi_k(\vek x) p_k=\col\psi^T(\vek x)\col{ p} \\
  \ten c_h&=\tsum_k \theta_k(\vek x)\ten c_k=\col\theta^T(\vek x)\col{\ten c}
  \label{eq:thetaDG}
\end{align}
where $\col\phi(\vek x)$, $\col\psi(\vek x)$, $\col\theta(\vek x)$
are global shape functions. Note, that most terms in the momentum balance and
the continuity equation can be build using standard (Navier-) Stokes elements.

We will be using the DEVSS formulation
(see Sec.~\ref{sec:weakviscoelastic}) and the
log-conformation scheme can be used as an option.

\subsection{Time discretization (without inertia)}
\label{sec:timeDG}

The weak form of the constitutive equation \eqref{eq:weakcDG}
can be written as
\begin{equation}
   \mat M\utilde{\dot c}=\utilde f(\utilde c,\vek u,t)
\end{equation}
where $\utilde c$ is the vector of conformation unknowns and $\mat M$ is a mass
matrix of the form
\begin{equation}
   \mat M=\begin{pmatrix}
              \mat M_1 & \mat 0 & \hdotsfor{2} & \mat 0 \\
              \mat 0 & \mat M_2 & \mat 0 & \ldots & \mat 0 \\
                              \vdots &  && \ddots  & \vdots \\
              \mat 0 & \hdotsfor{2} & \mat 0 & \mat M_{N_e}
      
          \end{pmatrix}
\end{equation}
with $\mat M_e$ the mass matrix of element $e$. The matrix $\mat M$ is constant
for a fixed Eulerian mesh and can easily be inverted:
\begin{equation}
   \mat M^{-1}=\begin{pmatrix}
              \mat M_1^{-1} & \mat 0 & \hdotsfor{2} & \mat 0 \\
              \mat 0 & \mat M_2^{-1} & \mat 0 & \ldots & \mat 0 \\
                              \vdots &  && \ddots  & \vdots \\
              \mat 0 & \hdotsfor{2} & \mat 0 & \mat M_{N_e}^{-1}
      
          \end{pmatrix}
\end{equation}
and it is possible to apply the inversion at element level. It is even possible
to choose the shape function based on unknowns in Gauss points, 
%as in done in
%\dyn\ \cite{Hulsen05b}, 
in which case the mass matrix is a diagonal matrix.	

Once a routine exists for computing $\utilde f(\utilde c,\vek u,t)$ any
standard \emph{explicit} time integration scheme can be easily implemented,
for example Euler forward:
\begin{equation}
   \mat M\dfrac{\utilde c^{n+1}-\utilde c^n}{\Delta t}
                               =\utilde f(\utilde c^n,\vek u^n,t^n)
\end{equation}
or second-order Adams-Bashforth (AB2):
\begin{equation}
   \mat M\dfrac{\utilde c^{n+1}-\utilde c^n}{\Delta t}
                               =\frac32\utilde f(\utilde c^n,\vek u^n,t^n)
                  -\frac12\utilde f(\utilde c^{n-1},\vek u^{n-1},t^{n-1})
\end{equation}

Remarks:
\begin{enumerate}
   \item In the examples below we build $\mat M$ as a system matrix and use the
         default solver to perform an LU-decomposition only once.
   \item Inflow boundary conditions at $\Gamma_{\text{in}}$ (see
         Eq.~\eqref{eq:cinflow}) cannot be imposed as Dirichlet boundary
         condition. The value of $\ten c_\text{in}$ must be given as a
         function when $\vek u\cdot\vek n<0$ on the boundary.
   \item For shape functions $\theta_k$ of first-order ($P_1^d$ or $Q_1^d$) 
         on element level the forward Euler scheme is unstable for \emph{any}
         time step $\Delta t$ with DG discretization of the convection terms
         (see \cite{Hulsen91} and the references therein).
         The relaxation terms can make it stable again, but time steps can be
         very limited. Therefore it is recommended to use AB2 for stability,
         which gives a proper CFL condition for the time step \cite{Hulsen01}.
         Other second-order explicit schemes, such a RK2, are also possible
         \cite{Hulsen91}.
         Similarly a time integration scheme of order $k+1$ is need for
         $k$-th order interpolation in space.
\end{enumerate}

After we have computed $\ten c^{n+1}$ we compute $\vek u^{n+1}$, 
$\vek p^{n+1}$ and $\ten E^{n+1}$ from the momentum balance,
 continuity equation and projection equation at time $t^{n+1}$:
\begin{align}
      -\nabla\cdot\big(2\eta_s\ten D(\vek u^{n+1})\big) +\nabla p^{n+1} 
  -2\alpha\nabla\cdot\big( \ten D(\vek u^{n+1}) -\ten E^{n+1} \big)
&=\nabla\cdot\ten \tau(\ten c^{n+1})
        + \rho\vek b \label{eq:mombalDG}\\
     \nabla\cdot\vek u^{n+1}&=0\\
   -\ten D(\vek u^{n+1})+\ten E^{n+1} &= \ten 0
\end{align}
We call this scheme the \emph{explicit stress} formulation: the stress tensor
$\ten \tau(\ten c^{n+1})$ in the momentum balance is computed in an
explicit way from $\ten c^{n+1}$, which itself has been computed first from the
constitutive equation.

The time integration procedures described above are only suitable for fluids
where $\eta_s\neq0$. For fluids with $\eta_s=0$, such as melts described by a
large number of modes, other procedures are required. For small values of
$\eta_s$ the time step limitation can be very severe and the explicit stress
formulation cannot be used either. 

\subsection{Planar Couette flow of an Oldroyd-B fluid for a
 Weissenberg number (Wi) of 3 and 5}

We consider an Oldroyd fluid for a planar Couette problem on a unit square.
The problem is similar to the problem for an Oldroyd-B fluid in
Sec.~\ref{sec:couette}, but now using the DG method with explicit
time-integration.
There are two problems \texttt{couette7.f90} (first-order time
integration) and \texttt{couette8.f90} (second-order time integration AB2).
For each main file there are two input files:
\begin{quote}
   \texttt{input\_couette7a.txt}: $\Wi=3$, time step $\Delta t=5\times10^{-3}$\\
   \texttt{input\_couette7b.txt}: $\Wi=3$, time step $\Delta t=2\times10^{-3}$\\
   \texttt{input\_couette8a.txt}: $\Wi=3$, time step $\Delta t=1\times10^{-2}$\\
   \texttt{input\_couette8b.txt}: $\Wi=5$, time step $\Delta t=1\times10^{-2}$
\end{quote}
It is clear by running the examples that:
\begin{itemize}
 \item for the Couette problem time steps for the
   first-order scheme need to be much smaller than for the second-order scheme
to get similar results. The first example \texttt{input\_couette7a.txt} even is
unstable.
  \item the Couette problem is unstable for $\Wi=5$, which is one of the main
drawbacks of DEVSS/DG compared to DEVSS-G/SUPG. It should be noted here that a
square mesh is a bad case: stability improves when elements are stretched in
$x$-direction. In many practical problems elements have been stretched in flow
direction.
\end{itemize}

\subsection{Axi-symmetric problem: flow around a sphere in
 cylindrical tube filled with a Giesekus fluid.}

We consider the problem of a sphere moving along
the center line of a cylindrical tube with a constant velocity of $U$.
The problem is similar to the problem in Sec.\ref{sec:sphere2}, but instead of
the DEVSS-G/SUPG method we use the DEVSS/DG method here.
For the viscoelastic problem we have the additional boundary
condition that it is assumed that the stress is zero at the inflow boundary.

There are two example files: 
\texttt{sphere4.f90} (first-order integration) and
\texttt{sphere5.f90} (second-order integration AB2).
The input files are \texttt{input\_sphere4.txt} and \texttt{input\_sphere5.txt}, respectively.
For this problem the time step limitations of the first-order and second-order
scheme are similar.
In contrast to the DEVSS-G/SUPG there is no need for a restart from a lower
Weissenberg number solution.
The post-processing files are \texttt{sphere\_post.f90} and 
\texttt{streamfunction.f90}.


\subsection{Flow around a cylinder confined between two walls for a FENE-P
model}
\label{sec:cylinderfenepdg}

This problem is a modification of the 
2D flow around a cylinder confined between two walls using a FENE-P model,
described in Sec.~\ref{sec:cylinderfenep} It uses small time steps to avoid the
trace of the conformation tensor becoming too large.

The example file for the 2D problem is \texttt{cylinder2.f90} and
and can be obtained by typing at the
command line:
\begin{quote}
   \texttt{getexample confined\_cylinder\_FENE-P}
\end{quote}
The mesh is generated by \tfem\ in the
following way:
\begin{quote}
   \texttt{mesh\_confined\_cylinder < meshinput1.txt}
\end{quote}

There is a main program \texttt{streamfunction.f90} to
generate the stream function. 
The program \texttt{cylinder\_post.f90}
creates post-processing data for plotting.

